\uuid{imVJ}
\exo7id{3736}
\titre{exo7 3736}
\auteur{quercia}
\organisation{exo7}
\datecreate{2010-03-11}
\isIndication{false}
\isCorrection{false}
\chapitre{Espace euclidien, espace normé}
\sousChapitre{Projection, symétrie}
\cpge{mpsi-mp2i.espaces-vectoriels.d, pcsi.espaces-vectoriels.d, ptsi.espaces-vectoriels.d, mpsi-mp2i.espaces-prehilbertiens.c, pcsi.espaces-prehilbertiens.c, mpsi-mp2i.espaces-prehilbertiens.d, pcsi.espaces-prehilbertiens.d, mp-mpi.endomorphismes-euclidien.c, pc.endomorphismes-euclidiens.a, psi.endomorphismes-euclidien.a, pt.isometries-euclidien.a, mp-mpi.endomorphismes-euclidien.f, pc.endomorphismes-euclidiens.d, psi.endomorphismes-euclidien.f, pt.isometries-euclidien.e}
\module{Algèbre}
\niveau{L2}
\difficulte{}

\contenu{
\question{
Soit $E$ un espace euclidien de dimension $n$.
Soient $F,G$ deux sous-espaces de $E$ tels que $F \perp G$. On note $s_F$
et $s_G$ les symétries orthogonales de bases $F$ et~$G$.
Montrer que $s_F\circ s_G = s_G \circ s_F = s_{(F\oplus G)^\perp}$.
}
}
